% =====================================================================
%  Chapitre 0 -- Les outils mathématiques du cours
%  Introduction au Machine Learning -- CNFSDP -- L3 S5 -- 2026-2027
%  Version niveau licence (12/10/2026)
%  Compilation : pdflatex chapitre0.tex (deux fois)
% =====================================================================
\documentclass[11pt]{article}
\usepackage{cnfsdp-poly}

\chapitre{0}{Les outils mathématiques du cours}
         {Matrices, minimisation, probabilités, conditionnement, estimation}

\begin{document}
\pagedetitre

\begin{objectifs}
Ce chapitre n'a pas de séance : il se lit chez soi, crayon en main, avant le chapitre~1.
Il rassemble \emph{uniquement} les outils mathématiques utilisés dans les chapitres 1 à 5, avec pour chacun l'endroit où il sert.
\begin{itemize}
  \item Vecteurs et matrices : tableau de données, matrice de covariance, valeurs propres (chapitres 2 et 4).
  \item Minimiser une fonction : dérivée, gradient, descente de gradient (chapitres 1 à 3).
  \item Probabilités : espérance, variance, loi normale, loi des grands nombres (chapitre 1).
  \item Probabilités conditionnelles : formule de Bayes, moyenne conditionnelle (chapitres 1 et 3).
  \item Estimation : biais, variance, intervalle de confiance, maximum de vraisemblance (chapitres 1 à 3).
  \item Outils informatiques : installation de Python, VS Code, Git (TP0).
\end{itemize}
Chaque section se termine par des exercices ; les corrigés sont regroupés à la fin du chapitre.
\end{objectifs}

\tableofcontents
\newpage

\section*{Notations}
\addcontentsline{toc}{section}{Notations}

\begin{center}
\renewcommand{\arraystretch}{1.25}
\begin{tabular}{>{\raggedright}p{0.3\textwidth} p{0.6\textwidth}}
\toprule
\textbf{Notation} & \textbf{Signification} \\
\midrule
$x\in\R^d$ & vecteur (en colonne) de $d$ nombres $x_1,\dots,x_d$ \\
$A\T$ & transposée de la matrice $A$ (lignes et colonnes échangées) \\
$\sum_{i=1}^n a_i$ & somme $a_1+a_2+\dots+a_n$ \\
$\bar y=\frac1n\sum_{i=1}^ny_i$ & moyenne des $n$ nombres $y_1,\dots,y_n$ \\
$\P(A)$ & probabilité de l'événement $A$ \\
$\E[X]$, $\Var(X)$ & espérance (moyenne théorique) et variance de $X$ \\
$\P(A\mid B)$, $\E[Y\mid X=x]$ & probabilité de $A$ sachant $B$ ; moyenne de $Y$ sachant $X=x$ \\
$\mathcal N(\mu,\sigma^2)$ & loi normale de moyenne $\mu$ et de variance $\sigma^2$ \\
$\Phi$ & fonction de répartition de la loi $\mathcal N(0,1)$ \\
$\argmin_c g(c)$ & la (ou les) valeur(s) de $c$ où $g$ atteint son minimum \\
\bottomrule
\end{tabular}
\end{center}

Les variables aléatoires sont notées en majuscules ($X$, $Y$), les valeurs observées en minuscules ($x$, $y$).

\section{Vecteurs et matrices}
\label{sec:algebre}

Un jeu de données est un tableau de nombres. Les vecteurs et les matrices permettent de le manipuler d'un bloc ; ils servent surtout à la régression linéaire (chapitre~2) et à l'ACP (chapitre~4).

\subsection{Vecteurs}

\begin{definition}[Vecteur, produit scalaire, norme]
Un \emph{vecteur} de $\R^d$ est une liste de $d$ nombres $x=(x_1,\dots,x_d)$, écrite en colonne.
Le \emph{produit scalaire} de deux vecteurs $x$ et $y$ est le nombre $\ps{x}{y}=x_1y_1+\dots+x_dy_d$.
La \emph{norme} (longueur) de $x$ est $\norm{x}=\sqrt{\ps{x}{x}}=\sqrt{x_1^2+\dots+x_d^2}$.
Deux vecteurs sont \emph{orthogonaux} si $\ps{x}{y}=0$.
\end{definition}

\begin{exemple}
Pour $x=(1,2,2)$ et $y=(2,0,-1)$ : $\ps{x}{y}=2+0-2=0$, donc $x$ et $y$ sont orthogonaux, et $\norm{x}=\sqrt{1+4+4}=3$.
Un vecteur de norme $1$ est dit \emph{unitaire} : il indique une direction.
\end{exemple}

\subsection{Matrices}

\begin{definition}[Matrice, transposée, produit]
Une \emph{matrice} $A$ de taille $n\times d$ est un tableau de nombres à $n$ lignes et $d$ colonnes ; $A_{ij}$ est le nombre en ligne $i$, colonne $j$.
Sa \emph{transposée} $A\T$, de taille $d\times n$, échange lignes et colonnes : $(A\T)_{ij}=A_{ji}$.
Si $A$ est de taille $n\times d$ et $B$ de taille $d\times p$, le \emph{produit} $AB$ est la matrice $n\times p$ de coefficients
\[
  (AB)_{ij}=\sum_{k=1}^dA_{ik}B_{kj}\qquad\text{(ligne $i$ de $A$ « fois » colonne $j$ de $B$).}
\]
\end{definition}

\begin{exemple}
$\begin{pmatrix}1&2\\3&4\end{pmatrix}\begin{pmatrix}5\\6\end{pmatrix}=\begin{pmatrix}1\times5+2\times6\\3\times5+4\times6\end{pmatrix}=\begin{pmatrix}17\\39\end{pmatrix}$.
Attention : en général $AB\ne BA$.
\end{exemple}

\begin{definition}[Matrice identité, inverse]
La matrice \emph{identité} $I_d$ a des $1$ sur la diagonale et des $0$ ailleurs : $I_dx=x$.
Une matrice carrée $A$ est \emph{inversible} s'il existe une matrice $A^{-1}$ telle que $AA^{-1}=A^{-1}A=I_d$.
Le système linéaire $A\beta=b$ a alors une unique solution, $\beta=A^{-1}b$.
\end{definition}

\begin{proposition}[Inverse d'une matrice $2\times2$]
$A=\begin{pmatrix}a&b\\c&d\end{pmatrix}$ est inversible si et seulement si $ad-bc\ne0$, et alors
$A^{-1}=\dfrac{1}{ad-bc}\begin{pmatrix}d&-b\\-c&a\end{pmatrix}$.
\end{proposition}
\begin{preuve}
Un calcul direct donne $\begin{pmatrix}a&b\\c&d\end{pmatrix}\begin{pmatrix}d&-b\\-c&a\end{pmatrix}=(ad-bc)\,I_2$ : si $ad-bc\ne0$, on divise par $ad-bc$.
(La réciproque, si $ad-bc=0$, est admise.)
\end{preuve}

\subsection{Le tableau de données et la matrice de covariance}

\begin{definition}[Tableau de données, centrage, covariance empirique]
On range les observations dans une matrice $X$ de taille $n\times d$ : la ligne $i$ décrit l'individu $i$, la colonne $j$ la variable $j$.
Le \emph{vecteur des moyennes} est $\bar x=(\bar x_1,\dots,\bar x_d)$, où $\bar x_j$ est la moyenne de la colonne $j$.
Le tableau \emph{centré} $X_c$ s'obtient en retranchant à chaque colonne sa moyenne.
La \emph{matrice de covariance empirique} est la matrice $d\times d$
\[
  S=\frac1n\,X_c\T X_c ,
\]
dont le coefficient $S_{jk}$ est la covariance empirique des variables $j$ et $k$, et le coefficient diagonal $S_{jj}$ la variance empirique de la variable $j$.
\end{definition}

\begin{exemple}
Trois individus, deux variables : $X=\begin{pmatrix}1&2\\2&2\\3&5\end{pmatrix}$, de moyennes $\bar x=(2,3)$.
Alors $X_c=\begin{pmatrix}-1&-1\\0&-1\\1&2\end{pmatrix}$ et
$S=\frac13\begin{pmatrix}2&3\\3&6\end{pmatrix}=\begin{pmatrix}2/3&1\\1&2\end{pmatrix}$.
La covariance positive ($S_{12}=1$) indique que les deux variables augmentent ensemble.
\end{exemple}

La matrice $S$ est \emph{symétrique} : $S\T=S$.

\subsection{Valeurs propres et vecteurs propres}

\begin{definition}[Valeur propre, vecteur propre]
Soit $A$ une matrice carrée. Un vecteur $u\ne0$ est un \emph{vecteur propre} de $A$, associé à la \emph{valeur propre} $\lambda$, si
$Au=\lambda u$ : la matrice ne fait qu'allonger ou raccourcir $u$, sans changer sa direction.
\end{definition}

\begin{theoreme}[Matrices symétriques, admis]
\label{thm:spectral}
Une matrice symétrique $A$ de taille $d\times d$ possède $d$ valeurs propres réelles $\lambda_1\ge\lambda_2\ge\dots\ge\lambda_d$, et l'on peut choisir des vecteurs propres associés $u_1,\dots,u_d$ unitaires et deux à deux orthogonaux.
Si $A$ est une matrice de covariance, toutes ses valeurs propres sont positives ou nulles.
\end{theoreme}

\begin{exemple}
Pour $A=\begin{pmatrix}2&1\\1&2\end{pmatrix}$ : $A\begin{pmatrix}1\\1\end{pmatrix}=\begin{pmatrix}3\\3\end{pmatrix}$ et $A\begin{pmatrix}1\\-1\end{pmatrix}=\begin{pmatrix}1\\-1\end{pmatrix}$.
Les valeurs propres sont $3$ et $1$, de vecteurs propres unitaires $u_1=\frac{1}{\sqrt2}(1,1)$ et $u_2=\frac{1}{\sqrt2}(1,-1)$, qui sont bien orthogonaux.
\end{exemple}

\begin{remarque}[À quoi cela sert]
En ACP (chapitre~4), on cherche la direction $u$ (vecteur unitaire) selon laquelle les individus sont le plus dispersés.
La variance des individus projetés sur $u$ vaut $u\T Su$ ; elle est maximale pour le premier vecteur propre $u_1$ de $S$, et vaut alors $\lambda_1$.
\end{remarque}

\subsection*{Exercices}

\begin{exercice}[Produit de matrices]
\label{ex:produit}
Calculer $A\T A$ pour $A=\begin{pmatrix}1&0\\1&1\\1&2\end{pmatrix}$, puis l'inverse de $A\T A$.
\end{exercice}

\begin{exercice}[Covariance empirique]
\label{ex:covariance}
Quatre ménages ont pour taille $2,4,4,6$ et pour nombre de pièces $1,2,3,2$. Calculer la matrice de covariance empirique $S$ de ces deux variables, puis la corrélation $S_{12}/\sqrt{S_{11}S_{22}}$.
\end{exercice}

\begin{exercice}[Valeurs propres]
\label{ex:spectre}
Vérifier que $(1,1)$ et $(1,-1)$ sont des vecteurs propres de $B=\begin{pmatrix}5&3\\3&5\end{pmatrix}$ et donner les valeurs propres associées.
\end{exercice}

\section{Minimiser une fonction}
\label{sec:optimisation}

Apprendre, c'est presque toujours chercher les paramètres qui rendent une erreur la plus petite possible. Il faut donc savoir trouver le minimum d'une fonction.

\subsection{Fonctions d'une variable}

\begin{definition}[Minimum, minimiseur]
\label{def:minimum}
On dit qu'une fonction $g$ \emph{atteint son minimum} en $c^*$ si $g(c^*)\le g(c)$ pour tout $c$.
Le point $c^*$ est alors un \emph{minimiseur} de $g$, et $g(c^*)$ la \emph{valeur minimale}. On note $c^*\in\argmin_cg(c)$.
\end{definition}

\begin{proposition}[Rappels, admis]
\label{prop:min1d}
Soit $g$ une fonction deux fois dérivable sur $\R$.
\begin{enumerate}[label=(\roman*)]
  \item Si $g$ atteint son minimum en $c^*$, alors $g'(c^*)=0$.
  \item Si $g''(c)\ge0$ pour tout $c$ (on dit que $g$ est \emph{convexe} : sa courbe est « en cuvette »), alors tout point où $g'$ s'annule est un minimiseur.
  \item Si $g''(c)>0$ pour tout $c$, il y a au plus un minimiseur. (C'est une condition suffisante, pas nécessaire : $c\mapsto c^4$ n'a qu'un minimiseur alors que sa dérivée seconde s'annule en $0$.)
\end{enumerate}
\end{proposition}

\begin{proposition}[La moyenne minimise les écarts au carré]
\label{prop:moyenne}
Pour des nombres $y_1,\dots,y_n$, la fonction $g(c)=\sum_{i=1}^n(y_i-c)^2$ atteint son minimum en un unique point, la moyenne $\bar y$.
\end{proposition}
\begin{preuve}
$g'(c)=-2\sum_i(y_i-c)=-2\big(\sum_iy_i-nc\big)$ s'annule pour $c=\frac1n\sum_iy_i=\bar y$, et $g''(c)=2n>0$ : on conclut par la proposition~\ref{prop:min1d}.
\end{preuve}

\begin{remarque}
Si l'on minimise plutôt $\sum_i|y_i-c|$ (écarts absolus), le minimiseur est une \emph{médiane} des $y_i$ (admis).
Ces deux résultats expliquent pourquoi la perte quadratique conduit à prédire des moyennes et la perte absolue des médianes (chapitre~1).
\end{remarque}

\subsection{Fonctions de plusieurs variables}

\begin{definition}[Dérivées partielles, gradient]
Pour une fonction $L(a,b)$ de deux variables, la \emph{dérivée partielle} $\frac{\partial L}{\partial a}$ est la dérivée par rapport à $a$, $b$ étant traité comme une constante (et de même pour $b$).
Le \emph{gradient} est le vecteur $\nabla L=\big(\frac{\partial L}{\partial a},\frac{\partial L}{\partial b}\big)$ ; il indique la direction dans laquelle $L$ augmente le plus vite.
\end{definition}

Comme en dimension $1$ : si $L$ atteint son minimum en un point, son gradient y est nul ; et si $L$ est convexe (en forme de cuvette), un point où le gradient s'annule est un minimiseur (admis).

\begin{proposition}[Droite des moindres carrés]
\label{prop:droite}
Soient des couples $(x_1,y_1),\dots,(x_n,y_n)$ avec des $x_i$ non tous égaux. La fonction
$L(a,b)=\sum_{i=1}^n\big(y_i-a-bx_i\big)^2$ atteint son minimum en
\[
  \hat b=\frac{s_{xy}}{s_x^2},\qquad \hat a=\bar y-\hat b\,\bar x,
\]
où $s_x^2=\frac1n\sum_i(x_i-\bar x)^2$ est la variance empirique des $x_i$ et $s_{xy}=\frac1n\sum_i(x_i-\bar x)(y_i-\bar y)$ la covariance empirique.
\end{proposition}
\begin{preuve}
$\frac{\partial L}{\partial a}=-2\sum_i(y_i-a-bx_i)=0$ donne $a=\bar y-b\bar x$.
En reportant, $\frac{\partial L}{\partial b}=-2\sum_ix_i\big((y_i-\bar y)-b(x_i-\bar x)\big)=0$. Comme $\sum_i\bar x\big((y_i-\bar y)-b(x_i-\bar x)\big)=0$, on peut remplacer $x_i$ par $x_i-\bar x$, ce qui donne
$\sum_i(x_i-\bar x)(y_i-\bar y)=b\sum_i(x_i-\bar x)^2$, soit $b=s_{xy}/s_x^2$. La fonction $L$ étant convexe (admis), ce point est le minimiseur.
\end{preuve}

\begin{exemple}
Pour $x=(1,2,3)$ et $y=(2,2,5)$ : $\bar x=2$, $\bar y=3$, $s_x^2=\frac23$, $s_{xy}=1$, donc $\hat b=1{,}5$ et $\hat a=0$. La droite $y=1{,}5\,x$ est celle qui passe « au plus près » des trois points.
C'est le point de départ de la régression linéaire du chapitre~2.
\end{exemple}

\subsection{La descente de gradient}

Quand on ne sait pas résoudre $\nabla L=0$ à la main (c'est le cas de la régression logistique, chapitre~3), on s'approche du minimum pas à pas.

\begin{definition}[Descente de gradient]
Partant d'un point $\theta_0$, on répète $\theta_{k+1}=\theta_k-\eta\,\nabla L(\theta_k)$, où $\eta>0$ est le \emph{pas} : on fait un petit pas dans la direction où $L$ descend le plus vite.
\end{definition}

\begin{proposition}[Un exemple où tout se calcule]
\label{prop:gd}
Pour $L(\theta)=(\theta-3)^2$, la descente de gradient de pas $\eta$ vérifie $\theta_k-3=(1-2\eta)^k(\theta_0-3)$.
Elle converge vers le minimiseur $3$ si et seulement si $0<\eta<1$.
\end{proposition}
\begin{preuve}
$L'(\theta)=2(\theta-3)$, donc $\theta_{k+1}-3=\theta_k-3-2\eta(\theta_k-3)=(1-2\eta)(\theta_k-3)$, et l'on itère.
La suite $(1-2\eta)^k$ tend vers $0$ si et seulement si $|1-2\eta|<1$, c'est-à-dire $0<\eta<1$.
\end{preuve}

\begin{remarque}[Lecture]
Un pas trop petit avance lentement ; un pas trop grand fait « sauter » par-dessus le minimum et peut diverger. Avec $\theta_0=0$ et $\eta=0{,}25$ : $\theta_1=1{,}5$, $\theta_2=2{,}25$, $\theta_3=2{,}625$\dots{} vers $3$.
\end{remarque}

\subsection*{Exercices}

\begin{exercice}[La meilleure constante]
\label{ex:constante}
Les revenus mensuels (en milliers de FCFA) de cinq ménages sont $80,100,120,150,550$.
Quelle constante minimise la somme des écarts au carré ? la somme des écarts absolus ? Commenter l'effet du revenu $550$.
\end{exercice}

\begin{exercice}[Droite des moindres carrés]
\label{ex:droite}
Calculer la droite des moindres carrés pour les points $(0,1),(1,3),(2,4),(3,8)$.
\end{exercice}

\begin{exercice}[Choix du pas]
\label{ex:gd}
Pour $L(\theta)=2\theta^2$, écrire la descente de gradient de pas $\eta$, puis trouver les valeurs de $\eta$ pour lesquelles elle converge vers $0$.
\end{exercice}

\section{Probabilités}
\label{sec:probabilites}

Les données sont vues comme les réalisations de variables aléatoires. C'est ce qui permet de parler d'« erreur moyenne » d'un modèle sur toute une population (chapitre~1).

\subsection{Variables aléatoires, espérance, variance}

\begin{definition}[Variable aléatoire, loi]
Une \emph{variable aléatoire} $X$ est une quantité dont la valeur dépend du hasard (l'individu tiré au sort, par exemple).
Elle est \emph{discrète} si elle prend des valeurs isolées $x_1,x_2,\dots$ avec des probabilités $p_k=\P(X=x_k)$ ;
elle est \emph{continue} si elle a une \emph{densité} $f\ge0$ telle que $\P(a\le X\le b)=\int_a^bf(x)\,dx$.
\end{definition}

\begin{definition}[Espérance, variance, écart-type]
L'\emph{espérance} (moyenne théorique) de $X$ est $\E[X]=\sum_kx_kp_k$ dans le cas discret, $\E[X]=\int xf(x)\,dx$ dans le cas continu.
La \emph{variance} mesure la dispersion autour de la moyenne : $\Var(X)=\E\big[(X-\E[X])^2\big]$ ; l'\emph{écart-type} est $\sqrt{\Var(X)}$.
\end{definition}

\begin{proposition}[Règles de calcul]
\label{prop:regles}
Pour des variables $X,Y$ et des réels $a,b$ :
\begin{enumerate}[label=(\roman*)]
  \item $\E[aX+bY]=a\E[X]+b\E[Y]$ (linéarité) ;
  \item $\Var(X)=\E[X^2]-(\E[X])^2$ et $\Var(aX+b)=a^2\Var(X)$.
\end{enumerate}
\end{proposition}
\begin{preuve}
(i) est admis. (ii) : en développant, $\E[(X-\E X)^2]=\E[X^2]-2\E[X]\E[X]+(\E X)^2$. Puis $aX+b-\E[aX+b]=a(X-\E X)$, dont le carré vaut $a^2(X-\E X)^2$.
\end{preuve}

\begin{proposition}[La meilleure constante est l'espérance]
\label{prop:meilleureconstante}
Pour toute variable $Z$ de variance finie et tout réel $c$ :
\[
  \E\big[(Z-c)^2\big]=\Var(Z)+\big(\E[Z]-c\big)^2 .
\]
Donc $c\mapsto\E[(Z-c)^2]$ atteint son minimum en $c=\E[Z]$, et la valeur minimale est $\Var(Z)$.
\end{proposition}
\begin{preuve}
On écrit $Z-c=(Z-\E Z)+(\E Z-c)$ et l'on développe :
$\E[(Z-c)^2]=\Var(Z)+2(\E Z-c)\E[Z-\E Z]+(\E Z-c)^2$, et $\E[Z-\E Z]=0$.
\end{preuve}
C'est la version « population » de la proposition~\ref{prop:moyenne}, et c'est elle qui donne le meilleur prédicteur au chapitre~1.

\subsection{Covariance, corrélation, indépendance}

\begin{definition}[Covariance, corrélation]
$\Cov(X,Y)=\E\big[(X-\E X)(Y-\E Y)\big]$ et, si les variances sont non nulles, la \emph{corrélation} est
$\rho(X,Y)=\dfrac{\Cov(X,Y)}{\sqrt{\Var(X)\Var(Y)}}$, toujours comprise entre $-1$ et $1$ (admis).
\end{definition}

\begin{definition}[Indépendance]
$X$ et $Y$ sont \emph{indépendantes} si connaître la valeur de l'une ne change pas la loi de l'autre : $\P(X\in A,\,Y\in B)=\P(X\in A)\,\P(Y\in B)$ pour tous ensembles $A$, $B$.
\end{definition}

\begin{proposition}[admise]
\label{prop:indep}
Si $X$ et $Y$ sont indépendantes, alors $\Cov(X,Y)=0$. Dans tous les cas, $\Var(X+Y)=\Var(X)+\Var(Y)+2\Cov(X,Y)$.
\end{proposition}

\begin{remarque}
Une covariance nulle n'implique pas l'indépendance : la corrélation ne mesure que la liaison \emph{linéaire}.
\end{remarque}

\subsection{Lois usuelles}

\begin{center}
\renewcommand{\arraystretch}{1.3}
\begin{tabular}{llll}
\toprule
\textbf{Loi} & \textbf{Description} & $\E[X]$ & $\Var(X)$ \\
\midrule
Bernoulli $\mathcal B(p)$ & $X=1$ avec probabilité $p$, $0$ sinon & $p$ & $p(1-p)$ \\
Binomiale $\mathcal B(n,p)$ & nombre de succès en $n$ essais indépendants & $np$ & $np(1-p)$ \\
Uniforme $\mathcal U[a,b]$ & densité constante $\frac{1}{b-a}$ sur $[a,b]$ & $\frac{a+b}{2}$ & $\frac{(b-a)^2}{12}$ \\
Normale $\mathcal N(\mu,\sigma^2)$ & densité $\frac{1}{\sigma\sqrt{2\pi}}e^{-(x-\mu)^2/(2\sigma^2)}$ & $\mu$ & $\sigma^2$ \\
\bottomrule
\end{tabular}
\end{center}

Si $Z\sim\mathcal N(0,1)$, on note $\Phi(t)=\P(Z\le t)$ ; on retiendra $\P(-1{,}96\le Z\le1{,}96)\approx0{,}95$.

\subsection{Moyenne d'un échantillon et théorèmes limites}

\begin{proposition}[Espérance et variance d'une moyenne]
\label{prop:moyennevar}
Soient $X_1,\dots,X_n$ indépendantes et de même loi, de moyenne $\mu$ et de variance $\sigma^2$, et $\bar X_n=\frac1n\sum_iX_i$. Alors
$\E[\bar X_n]=\mu$ et $\Var(\bar X_n)=\dfrac{\sigma^2}{n}$.
\end{proposition}
\begin{preuve}
Par linéarité, $\E[\bar X_n]=\frac1n\cdot n\mu=\mu$. Par indépendance (proposition~\ref{prop:indep}), les covariances sont nulles et
$\Var(\bar X_n)=\frac{1}{n^2}\sum_i\Var(X_i)=\frac{n\sigma^2}{n^2}=\frac{\sigma^2}{n}$.
\end{preuve}

\begin{theoreme}[Loi des grands nombres, admis]
\label{thm:lgn}
Sous les mêmes hypothèses, $\bar X_n$ se rapproche de $\mu$ quand $n$ devient grand : pour tout $\varepsilon>0$, $\P(|\bar X_n-\mu|>\varepsilon)\to0$.
\end{theoreme}

\begin{theoreme}[Théorème central limite, admis]
\label{thm:tcl}
Pour $n$ grand, $\bar X_n$ suit approximativement la loi $\mathcal N\big(\mu,\sigma^2/n\big)$. En particulier,
$\P\big(|\bar X_n-\mu|\le1{,}96\,\sigma/\sqrt n\big)\approx0{,}95$.
\end{theoreme}

\begin{aretenir}
Une moyenne calculée sur $n$ observations se trompe typiquement de $\sigma/\sqrt n$ : pour diviser l'erreur par $2$, il faut $4$ fois plus de données.
C'est ce qui permet d'estimer l'erreur d'un modèle sur un échantillon test (chapitre~1).
\end{aretenir}

\subsection*{Exercices}

\begin{exercice}[Variable de Bernoulli]
\label{ex:bernoulli}
Soit $X\sim\mathcal B(p)$. Calculer $\E[X]$, $\E[X^2]$ et retrouver $\Var(X)=p(1-p)$. Pour quelle valeur de $p$ la variance est-elle maximale ?
\end{exercice}

\begin{exercice}[Précision d'un sondage]
\label{ex:sondage}
On estime la proportion $p$ de ménages équipés d'un réfrigérateur par la proportion observée sur $n$ ménages.
En utilisant le théorème central limite et $p(1-p)\le\frac14$, quelle taille $n$ garantit une erreur d'au plus $3$ points avec probabilité $95\,\%$ ?
\end{exercice}

\begin{exercice}[Meilleure constante]
\label{ex:meilleureconstante}
Un dé équilibré donne $Z\in\{1,\dots,6\}$. Calculer $\E[(Z-c)^2]$ pour $c=3$ et pour $c=3{,}5$, et vérifier la proposition~\ref{prop:meilleureconstante}.
\end{exercice}

\section{Probabilités conditionnelles}
\label{sec:conditionnement}

Prédire $Y$ à partir de $X$, c'est raisonner « sachant $X$ ». Cette section donne les deux outils du chapitre~1 : la formule de Bayes et la moyenne conditionnelle.

\subsection{Probabilité conditionnelle et formule de Bayes}

\begin{definition}[Probabilité conditionnelle]
Pour deux événements $A$ et $B$ avec $\P(B)>0$, la probabilité de $A$ \emph{sachant} $B$ est $\P(A\mid B)=\dfrac{\P(A\text{ et }B)}{\P(B)}$.
\end{definition}

\begin{proposition}[Probabilités totales et formule de Bayes]
\label{prop:bayes}
Si $B_1,\dots,B_K$ forment une partition (un et un seul se produit), de probabilités non nulles :
\[
  \P(A)=\sum_{k=1}^K\P(A\mid B_k)\,\P(B_k)
  \qquad\text{et}\qquad
  \P(B_j\mid A)=\frac{\P(A\mid B_j)\,\P(B_j)}{\sum_k\P(A\mid B_k)\,\P(B_k)} .
\]
\end{proposition}
\begin{preuve}
$A$ est la réunion des événements disjoints « $A$ et $B_k$ », de probabilités $\P(A\mid B_k)\P(B_k)$ : on additionne.
Puis $\P(B_j\mid A)=\P(A\text{ et }B_j)/\P(A)$, et l'on remplace $\P(A)$ par la première formule.
\end{preuve}

\begin{exemple}[Test de dépistage]
Une maladie touche $1\,\%$ d'une population. Un test est positif chez $95\,\%$ des malades et négatif chez $90\,\%$ des non-malades. Pour une personne testée positive :
\[
  \P(M\mid+)=\frac{0{,}95\times0{,}01}{0{,}95\times0{,}01+0{,}10\times0{,}99}=\frac{0{,}0095}{0{,}1085}\approx0{,}088 .
\]
Moins de $9\,\%$ des positifs sont malades : la rareté de la maladie pèse lourd. C'est le raisonnement du classifieur de Bayes (chapitre~1).
\end{exemple}

\subsection{Moyenne conditionnelle}

\begin{definition}[Loi et espérance conditionnelles, cas discret]
Si $X$ et $Y$ sont discrètes et $\P(X=x)>0$, la \emph{loi de $Y$ sachant $X=x$} est donnée par les probabilités $\P(Y=y\mid X=x)$, et
\[
  \E[Y\mid X=x]=\sum_yy\;\P(Y=y\mid X=x)
\]
est la \emph{moyenne de $Y$ sachant $X=x$} : la moyenne de $Y$ dans le groupe des individus pour lesquels $X=x$.
\end{definition}

Quand $Y$ ne prend que les valeurs $0$ et $1$, $\E[Y\mid X=x]=\P(Y=1\mid X=x)$ : une moyenne de $0$ et de $1$ est une proportion.
Le cas continu se définit de la même façon avec des densités ; nous l'admettrons.

\begin{proposition}[Formule de l'espérance totale]
\label{prop:tour}
$\E[Y]=\sum_x\E[Y\mid X=x]\;\P(X=x)$ : la moyenne générale est la moyenne pondérée des moyennes de groupes.
\end{proposition}
\begin{preuve}
$\sum_x\E[Y\mid X=x]\P(X=x)=\sum_x\sum_yy\,\P(Y=y\mid X=x)\P(X=x)=\sum_y y\sum_x\P(Y=y,X=x)=\sum_yy\,\P(Y=y)=\E[Y]$.
\end{preuve}

\begin{exemple}
Parmi les femmes d'une région, $40\,\%$ vivent en ville, où la moyenne est de $2{,}5$ enfants, et $60\,\%$ en milieu rural, avec $4{,}5$ enfants en moyenne.
La moyenne régionale est $0{,}4\times2{,}5+0{,}6\times4{,}5=3{,}7$ enfants.
Pour prédire le nombre d'enfants d'une femme dont on connaît le milieu, on utilisera $2{,}5$ ou $4{,}5$ plutôt que $3{,}7$ (chapitre~1).
\end{exemple}

\subsection*{Exercices}

\begin{exercice}[Prévalence et dépistage]
\label{ex:prevalence}
Reprendre le test de dépistage avec une prévalence de $10\,\%$. Que devient $\P(M\mid+)$ ? Commenter.
\end{exercice}

\begin{exercice}[Moyennes conditionnelles]
\label{ex:tableau}
Le couple $(X,Y)$ a la loi suivante.
\begin{center}
\begin{tabular}{c|ccc}
 & $Y=0$ & $Y=1$ & $Y=2$ \\ \hline
$X=0$ & $0{,}2$ & $0{,}1$ & $0{,}1$ \\
$X=1$ & $0{,}1$ & $0{,}2$ & $0{,}3$
\end{tabular}
\end{center}
Calculer $\E[Y\mid X=0]$ et $\E[Y\mid X=1]$, puis vérifier la formule de l'espérance totale.
\end{exercice}

\section{Estimation}
\label{sec:statistique}

Un modèle de machine learning est fabriqué à partir d'un échantillon : il en hérite les erreurs d'estimation. Le vocabulaire de cette section (biais, variance) sert dès le chapitre~1.

\subsection{Biais, variance, erreur quadratique}

\begin{definition}[Estimateur, biais, erreur quadratique moyenne]
On cherche une quantité inconnue $\theta$ (une moyenne, une proportion\dots). Un \emph{estimateur} $\hat\theta$ est une formule calculée sur l'échantillon ; il est aléatoire, puisque l'échantillon l'est.
Son \emph{biais} est $\E[\hat\theta]-\theta$ (l'erreur systématique), et son \emph{erreur quadratique moyenne} est $\mathrm{EQM}(\hat\theta)=\E\big[(\hat\theta-\theta)^2\big]$.
\end{definition}

\begin{proposition}[Décomposition biais-variance]
\label{prop:biaisvariance}
$\mathrm{EQM}(\hat\theta)=\big(\E[\hat\theta]-\theta\big)^2+\Var(\hat\theta)$.
\end{proposition}
\begin{preuve}
C'est la proposition~\ref{prop:meilleureconstante} appliquée à $Z=\hat\theta$ et $c=\theta$.
\end{preuve}

\begin{remarque}[Lecture]
Une erreur vient soit d'un tir mal centré (biais), soit d'un tir dispersé (variance). Au chapitre~1, le même découpage explique le sur-apprentissage : un modèle trop simple est biaisé, un modèle trop complexe a une forte variance.
\end{remarque}

\begin{exemple}[Moyenne empirique]
D'après la proposition~\ref{prop:moyennevar}, $\bar X_n$ est sans biais pour $\mu$ et $\mathrm{EQM}(\bar X_n)=\sigma^2/n$.
\end{exemple}

\subsection{Intervalle de confiance d'une proportion}

\begin{proposition}[admise, conséquence du théorème central limite]
\label{prop:ic}
Si $\hat p$ est la proportion observée d'un caractère sur $n$ individus tirés au hasard, alors pour $n$ grand,
\[
  \Big[\hat p-1{,}96\sqrt{\tfrac{\hat p(1-\hat p)}{n}}\ ;\ \hat p+1{,}96\sqrt{\tfrac{\hat p(1-\hat p)}{n}}\Big]
\]
contient la vraie proportion $p$ avec une probabilité d'environ $95\,\%$.
\end{proposition}
Au chapitre~1, $\hat p$ sera le taux d'erreur d'un classifieur mesuré sur un échantillon test.

\subsection{Maximum de vraisemblance}

\begin{definition}[Vraisemblance, estimateur du maximum de vraisemblance]
Si les observations $x_1,\dots,x_n$ sont indépendantes, de loi dépendant d'un paramètre $\theta$ (probabilité ou densité $p_\theta$), la \emph{vraisemblance} est
$\mathcal L(\theta)=p_\theta(x_1)\times\dots\times p_\theta(x_n)$ : la probabilité d'avoir observé ces données si le paramètre vaut $\theta$.
L'\emph{estimateur du maximum de vraisemblance} est la valeur de $\theta$ qui rend $\mathcal L(\theta)$, ou de façon équivalente $\log\mathcal L(\theta)$, la plus grande.
\end{definition}

\begin{exemple}[Proportion]
Pour des observations $x_i\in\{0,1\}$ de loi $\mathcal B(\theta)$, avec $s=\sum_ix_i$ succès ($0<s<n$) :
$\log\mathcal L(\theta)=s\log\theta+(n-s)\log(1-\theta)$, de dérivée $\frac s\theta-\frac{n-s}{1-\theta}$, nulle pour $\hat\theta=s/n$ : la proportion observée.
\end{exemple}

\begin{remarque}
Maximiser la vraisemblance revient à minimiser $-\log\mathcal L(\theta)$, une somme d'erreurs sur les observations. La régression logistique du chapitre~3 est construite ainsi.
\end{remarque}

\subsection*{Exercices}

\begin{exercice}[Un biais qui paie]
\label{ex:shrinkage}
Soient $X_1,\dots,X_n$ indépendantes, de moyenne $\mu\ne0$ et de variance $\sigma^2$. Pour $c\in\R$, on estime $\mu$ par $c\,\bar X_n$.
Calculer le biais, la variance, puis l'erreur quadratique moyenne de $c\,\bar X_n$. Montrer qu'elle est minimale pour $c^*=\dfrac{\mu^2}{\mu^2+\sigma^2/n}<1$.
Pourquoi ce $c^*$ n'est-il pas directement utilisable ?
\end{exercice}

\begin{exercice}[Intervalle de confiance]
\label{ex:ic}
Un classifieur se trompe sur $45$ des $300$ individus d'un échantillon test. Donner un intervalle de confiance à $95\,\%$ pour son taux d'erreur.
\end{exercice}

\section{En Python}
\label{sec:python}

Chaque outil de ce chapitre se vérifie en quelques lignes de NumPy. Les sorties indiquées en commentaire sont celles obtenues ; ces programmes sont regroupés dans le notebook \texttt{TP0.ipynb}.

\subsection{Matrice de covariance et valeurs propres}
\begin{python}
import numpy as np

# Matrice de covariance et valeurs propres (sections 1)
X = np.array([[1, 2], [2, 2], [3, 5]])        # 3 individus, 2 variables
Xc = X - X.mean(axis=0)                        # tableau centré
S = Xc.T @ Xc / len(X)                         # covariance empirique
print(S.round(3))                              # [[0.667 1.] [1. 2.]]
lam, U = np.linalg.eigh(S)                     # valeurs propres croissantes
print(lam.round(4))                            # [0.1315 2.5352]
print(U[:, -1].round(3))                       # premier axe de l'ACP
\end{python}
On retrouve la matrice $S$ de l'exemple de la section~\ref{sec:algebre}. Le vecteur propre de la plus grande valeur propre donne la direction de plus grande dispersion des trois points : c'est le premier axe de l'ACP.

\subsection{Droite des moindres carrés et descente de gradient}
\begin{python}
import numpy as np

# Droite des moindres carrés (section 2)
x = np.array([1, 2, 3]); y = np.array([2, 2, 5])
b = np.mean((x - x.mean()) * (y - y.mean())) / np.var(x)   # s_xy / s_x²
a = y.mean() - b * x.mean()
print(a, b)                                    # 0.0 1.5
print(np.polyfit(x, y, deg=1).round(4))        # [1.5 0.] : même droite

# Descente de gradient sur L(θ) = (θ - 3)²
theta, eta = 0.0, 0.25
for k in range(6):
    theta = theta - eta * 2 * (theta - 3)
    print(round(theta, 4), end=" ")            # 1.5 2.25 2.625 2.8125 ...
\end{python}
Les formules de la proposition~\ref{prop:droite} donnent la même droite que \texttt{np.polyfit}, et la descente de gradient s'approche de $3$ comme prévu par la proposition~\ref{prop:gd}.

\subsection{Loi des grands nombres et théorème central limite}
\begin{python}
import numpy as np
rng = np.random.default_rng(2)

# Loi des grands nombres et TCL (section 3)
p, n, R = 0.3, 400, 20_000                     # R sondages de n personnes
props = rng.binomial(n, p, size=R) / n         # R proportions observées
print(props.mean().round(4), props.var().round(6), p * (1 - p) / n)
ecart = 1.96 * np.sqrt(p * (1 - p) / n)
print(np.mean(np.abs(props - p) <= ecart).round(3))   # proche de 0.95
\end{python}
La variance des proportions observées vaut bien $p(1-p)/n$, et l'écart à $p$ reste sous $1{,}96\sqrt{p(1-p)/n}$ dans environ $95\,\%$ des sondages simulés ($94{,}3\,\%$ ici).

\subsection{Moyennes conditionnelles}
\begin{python}
import numpy as np, pandas as pd
rng = np.random.default_rng(3)

# Moyennes conditionnelles par groupe (section 4)
n = 10_000
milieu = rng.choice(["urbain", "rural"], size=n, p=[0.4, 0.6])
enfants = np.where(milieu == "urbain", rng.poisson(2.5, n), rng.poisson(4.5, n))
df = pd.DataFrame({"milieu": milieu, "enfants": enfants})

print(df.groupby("milieu")["enfants"].mean().round(2))   # ≈ 4.5 et 2.5
part = df["milieu"].value_counts(normalize=True)
moy = df.groupby("milieu")["enfants"].mean()
print(round((part * moy).sum(), 4), round(df["enfants"].mean(), 4))  # espérance totale
\end{python}
Les moyennes par groupe estiment $\E[Y\mid X=x]$, et leur moyenne pondérée redonne la moyenne générale (proposition~\ref{prop:tour}).

\section{TP0 \textendash{} Installation et premiers pas}
\label{sec:tp0}

\begin{tp}{TP0 \textendash{} à faire chez soi avant le chapitre 1}{environ 1 h 30 \textperiodcentered{} \texttt{\cnfDepot}}
Objectif : arriver au premier cours avec un environnement Python fonctionnel, le dépôt du cours sur sa machine, et le notebook \texttt{TP0.ipynb} exécuté.
En cas de blocage, Google Colab permet de faire tous les TP dans un navigateur, sans rien installer (étape~6).
\end{tp}

\subsection*{Étape 1 \textendash{} Python avec Miniconda}
Télécharger l'installateur Miniconda pour votre système depuis le site officiel d'Anaconda (\url{https://www.anaconda.com/docs/getting-started/miniconda/install}) et l'installer avec les options par défaut.
Sous Windows, ouvrir ensuite le programme \emph{Anaconda Prompt} ; sous macOS ou Linux, un terminal. Vérifier :
\begin{terminal}
conda --version
\end{terminal}

\subsection*{Étape 2 \textendash{} Environnement du cours}
Créer un environnement isolé contenant les bibliothèques du cours, puis l'activer :
\begin{terminal}
conda create -n ml-cnfsdp -c conda-forge python=3.12 numpy pandas matplotlib scikit-learn jupyter ipykernel pyreadstat
conda activate ml-cnfsdp
python -c "import numpy, sklearn; print(numpy.__version__, sklearn.__version__)"
\end{terminal}
La bibliothèque \texttt{pyreadstat} servira à lire les fichiers SPSS des enquêtes DHS au chapitre~5.

\subsection*{Étape 3 \textendash{} VS Code}
Installer Visual Studio Code (\url{https://code.visualstudio.com}), puis, dans l'onglet \emph{Extensions}, les extensions \emph{Python} et \emph{Jupyter} publiées par Microsoft.
Dans VS Code, \texttt{Ctrl+Maj+P}, puis \emph{Python: Select Interpreter}, et choisir l'environnement \texttt{ml-cnfsdp}.

\subsection*{Étape 4 \textendash{} Git et Git Bash}
Installer Git (\url{https://git-scm.com}) ; sous Windows, l'installation fournit le terminal \emph{Git Bash}. Puis, une seule fois :
\begin{terminal}
git config --global user.name "Prénom Nom"
git config --global user.email "adresse@exemple.com"
\end{terminal}

\subsection*{Étape 5 \textendash{} GitHub et dépôt du cours}
Créer un compte gratuit sur \url{https://github.com}. Dans Git Bash, cloner le dépôt du cours et l'ouvrir dans VS Code :
\begin{terminal}
git clone https://github.com/rdimbamba/ml-cnfsdp.git
cd ml-cnfsdp
code .
\end{terminal}
Pour récupérer les nouveaux supports au fil du semestre : \texttt{git pull} depuis ce dossier.

\subsection*{Étape 6 \textendash{} Exécuter le notebook}
Ouvrir \texttt{notebooks/TP0.ipynb} dans VS Code et exécuter toutes les cellules (\emph{Run All}).
Sans installation locale : ouvrir \url{https://colab.research.google.com}, onglet \emph{GitHub}, saisir \texttt{rdimbamba/ml-cnfsdp}, puis choisir \texttt{TP0.ipynb} ; ou cliquer sur le bouton \emph{Open in Colab} de la page du cours.

\subsection*{Liste de vérification}
\begin{itemize}[label=$\square$]
  \item \texttt{conda activate ml-cnfsdp} fonctionne et \texttt{import sklearn} ne renvoie pas d'erreur.
  \item VS Code exécute une cellule de notebook avec l'environnement \texttt{ml-cnfsdp}.
  \item \texttt{git -{}-version} répond, et le dossier \texttt{ml-cnfsdp} est cloné.
  \item Les quatre programmes de la section~\ref{sec:python} donnent les mêmes sorties.
\end{itemize}

\subsection*{Exercices du TP0}
Ils sont posés dans le notebook ; chacun vérifie numériquement un résultat du chapitre.
\begin{enumerate}
  \item \emph{Matrices.} Calculer avec NumPy $A\T A$ et son inverse pour la matrice de l'exercice~\ref{ex:produit}, et comparer avec votre calcul à la main.
  \item \emph{Bayes par simulation.} Simuler un million de personnes avec la prévalence et le test de l'exemple de la section~\ref{sec:conditionnement}, et estimer $\P(M\mid+)$.
  \item \emph{Choix du pas.} Pour $L(\theta)=2\theta^2$, lancer $50$ pas de descente de gradient depuis $\theta_0=1$ avec $\eta=0{,}4$, puis $\eta=0{,}6$ ; comparer avec l'exercice~\ref{ex:gd}.
  \item \emph{Un biais qui paie.} Pour $\mu=1$, $\sigma=3$, $n=10$, estimer par simulation l'erreur quadratique de $c\,\bar X_n$ pour $c\in[0{,}5;1{,}2]$ et comparer le minimum à la valeur $c^*$ de l'exercice~\ref{ex:shrinkage}.
\end{enumerate}

\section{Corrigés des exercices}
\label{sec:corriges}

\begin{corrige}
\textbf{Exercice~\ref{ex:produit}.}
$A\T A=\begin{pmatrix}1&1&1\\0&1&2\end{pmatrix}\begin{pmatrix}1&0\\1&1\\1&2\end{pmatrix}=\begin{pmatrix}3&3\\3&5\end{pmatrix}$.
Son déterminant vaut $3\times5-3\times3=6\ne0$, donc $(A\T A)^{-1}=\frac16\begin{pmatrix}5&-3\\-3&3\end{pmatrix}$.
\end{corrige}

\begin{corrige}
\textbf{Exercice~\ref{ex:covariance}.}
Moyennes : $4$ personnes et $2$ pièces. Écarts : $(-2,0,0,2)$ et $(-1,0,1,0)$.
$S_{11}=\frac{4+0+0+4}{4}=2$, $S_{22}=\frac{1+0+1+0}{4}=0{,}5$, $S_{12}=\frac{2+0+0+0}{4}=0{,}5$, donc
$S=\begin{pmatrix}2&0{,}5\\0{,}5&0{,}5\end{pmatrix}$ et la corrélation vaut $\frac{0{,}5}{\sqrt{2\times0{,}5}}=0{,}5$.
\end{corrige}

\begin{corrige}
\textbf{Exercice~\ref{ex:spectre}.}
$B\begin{pmatrix}1\\1\end{pmatrix}=\begin{pmatrix}8\\8\end{pmatrix}=8\begin{pmatrix}1\\1\end{pmatrix}$ et $B\begin{pmatrix}1\\-1\end{pmatrix}=\begin{pmatrix}2\\-2\end{pmatrix}=2\begin{pmatrix}1\\-1\end{pmatrix}$ : valeurs propres $8$ et $2$.
\end{corrige}

\begin{corrige}
\textbf{Exercice~\ref{ex:constante}.}
La somme des écarts au carré est minimale pour la moyenne, $\frac{80+100+120+150+550}{5}=200$ ; la somme des écarts absolus pour la médiane, $120$.
Le revenu $550$ tire la moyenne vers le haut : quatre ménages sur cinq gagnent moins que la moyenne. La médiane résiste aux valeurs extrêmes.
\end{corrige}

\begin{corrige}
\textbf{Exercice~\ref{ex:droite}.}
$\bar x=1{,}5$, $\bar y=4$ ; écarts $(-1{,}5;-0{,}5;0{,}5;1{,}5)$ et $(-3,-1,0,4)$.
$s_{xy}=\frac{4{,}5+0{,}5+0+6}{4}=2{,}75$ et $s_x^2=\frac{2{,}25+0{,}25+0{,}25+2{,}25}{4}=1{,}25$, d'où $\hat b=2{,}2$ et $\hat a=4-2{,}2\times1{,}5=0{,}7$ : la droite est $y=0{,}7+2{,}2\,x$.
\end{corrige}

\begin{corrige}
\textbf{Exercice~\ref{ex:gd}.}
$L'(\theta)=4\theta$, donc $\theta_{k+1}=\theta_k-4\eta\theta_k=(1-4\eta)\theta_k$ et $\theta_k=(1-4\eta)^k\theta_0$.
La descente converge vers $0$ si et seulement si $|1-4\eta|<1$, c'est-à-dire $0<\eta<\frac12$.
\end{corrige}

\begin{corrige}
\textbf{Exercice~\ref{ex:bernoulli}.}
$\E[X]=1\times p+0\times(1-p)=p$ et $\E[X^2]=p$ (car $X^2=X$), donc $\Var(X)=p-p^2=p(1-p)$, maximale pour $p=\frac12$ où elle vaut $\frac14$.
\end{corrige}

\begin{corrige}
\textbf{Exercice~\ref{ex:sondage}.}
Il faut $1{,}96\sqrt{p(1-p)/n}\le0{,}03$. Comme $p(1-p)\le\frac14$, il suffit que $1{,}96\times\frac{1}{2\sqrt n}\le0{,}03$, soit $n\ge\big(\frac{1{,}96}{0{,}06}\big)^2\approx1\,067{,}1$ : $n=1\,068$ ménages.
\end{corrige}

\begin{corrige}
\textbf{Exercice~\ref{ex:meilleureconstante}.}
$\E[Z]=3{,}5$ et $\Var(Z)=\frac{35}{12}\approx2{,}917$.
Pour $c=3$ : $\E[(Z-3)^2]=\frac{4+1+0+1+4+9}{6}=\frac{19}{6}\approx3{,}167=\frac{35}{12}+(3{,}5-3)^2$.
Pour $c=3{,}5$ : on obtient $\Var(Z)$, la valeur minimale.
\end{corrige}

\begin{corrige}
\textbf{Exercice~\ref{ex:prevalence}.}
$\P(+)=0{,}95\times0{,}10+0{,}10\times0{,}90=0{,}185$ et $\P(M\mid+)=\frac{0{,}095}{0{,}185}\approx0{,}51$.
Avec le même test, un résultat positif passe de $9\,\%$ à $51\,\%$ de chances d'être un vrai malade : la prévalence compte autant que la qualité du test.
\end{corrige}

\begin{corrige}
\textbf{Exercice~\ref{ex:tableau}.}
$\P(X=0)=0{,}4$ et $\P(X=1)=0{,}6$.
$\E[Y\mid X=0]=\frac{0\times0{,}2+1\times0{,}1+2\times0{,}1}{0{,}4}=0{,}75$ et $\E[Y\mid X=1]=\frac{0\times0{,}1+1\times0{,}2+2\times0{,}3}{0{,}6}=\frac43$.
Espérance totale : $0{,}4\times0{,}75+0{,}6\times\frac43=1{,}1$, et directement $\E[Y]=1\times0{,}3+2\times0{,}4=1{,}1$.
\end{corrige}

\begin{corrige}
\textbf{Exercice~\ref{ex:shrinkage}.}
Biais : $c\mu-\mu=(c-1)\mu$ ; variance : $c^2\sigma^2/n$. Par la proposition~\ref{prop:biaisvariance},
$\mathrm{EQM}(c)=(1-c)^2\mu^2+c^2\frac{\sigma^2}{n}$, dont la dérivée $-2(1-c)\mu^2+2c\frac{\sigma^2}{n}$ s'annule en $c^*=\frac{\mu^2}{\mu^2+\sigma^2/n}<1$ (minimum, car la dérivée seconde est positive).
Accepter un peu de biais réduit davantage la variance : c'est l'idée de la régression ridge (chapitre~2).
Mais $c^*$ dépend de $\mu$ et $\sigma$, inconnus : en pratique, on règle ce type de paramètre sur les données, par validation croisée (chapitre~1).
\end{corrige}

\begin{corrige}
\textbf{Exercice~\ref{ex:ic}.}
$\hat p=\frac{45}{300}=0{,}15$ et $1{,}96\sqrt{\frac{0{,}15\times0{,}85}{300}}\approx0{,}040$ : le taux d'erreur est compris entre $11\,\%$ et $19\,\%$ avec une confiance d'environ $95\,\%$.
\end{corrige}


\end{document}
